%% 48장 제약과 쌍대성
\chapter{제약과 쌍대성}
\label{ch:constraints-duality}

\abstract{변수가 마음대로 움직일 수 없는 최적화 문제는 흔하다. 확률은 음수일 수 없고, 비율의 합은 $1$이며, 예산에는 한도가 있다. 이 장에서는 제약이 있는 최적화의 바탕을 정리한다. 먼저 42장의 사영 정리로 Farkas 보조정리를 증명하고, 그것으로 선형 제약 아래 국소 최소점이 만족하는 Karush--Kuhn--Tucker 조건을 끌어낸다. 이어서 라그랑주 쌍대 문제를 세워 약한 쌍대성을 보이고, 선형계획에서는 강한 쌍대성과 상보 여유 조건이 성립함을 증명한다. 선형계획을 푸는 두 방법으로 꼭짓점을 따라가는 심플렉스 방법과 중심 경로를 따라가는 내부점법을 다루는데, Klee--Minty 입방체에서 Dantzig의 규칙이 $2^n-1$번 피벗하는 것과 무작위 문제에서 내부점법의 반복 수가 크기에 거의 무관한 것을 확인한다. 끝으로 통계에서 자주 만나는 상자 제약 최소제곱, 곧 음이 아닌 최소제곱(NNLS)과 상하한 최소제곱(BVLS)의 활성 집합 방법을 다룬다.}

\section{제약 최적화와 KKT 조건}
\label{sec:kkt}

\subsection{문제와 라그랑주 함수}

일반적인 제약 최적화 문제는
\[
\min_{x\in\mathbb R^n}f(x)\quad\text{조건: }c_i(x)=0\ (i\in\mathcal E),\qquad c_i(x)\ge0\ (i\in\mathcal I)
\]
의 꼴이다. 제약을 만족하는 점을 \term{가능해}{feasible point}라 하고, 가능해 $x$에서 $c_i(x)=0$인 부등식 제약을 \term{활성 제약}{active constraint}이라 한다. 등식 제약은 언제나 활성이다. 활성 제약의 첨자 집합을 $\mathcal A(x)$로 쓴다. 승수 $\lambda_i$를 붙인
\[
\mathcal L(x,\lambda)=f(x)-\sum_{i\in\mathcal E\cup\mathcal I}\lambda_ic_i(x)
\]
를 \term{라그랑주 함수}{Lagrangian}라 한다.

제약이 없으면 최소점에서 $\nabla f=0$이었다(42장). 제약이 있으면 기울기가 $0$일 필요는 없고, 가능한 방향으로 움직여 $f$를 줄일 수 없기만 하면 된다. 이것을 식으로 옮기려면 가능한 방향의 집합을 대수적으로 표현해야 하는데, 제약이 선형이면 그것이 정확히 된다. 그 바탕이 Farkas의 보조정리이다.

\subsection{Farkas의 보조정리}

벡터 $a_1,\dots,a_k\in\mathbb R^n$의 음이 아닌 일차 결합 전체 $K=\{\sum_j\mu_ja_j:\mu_j\ge0\}$를 이 벡터들이 생성하는 \term{원뿔}{cone}이라 한다.

\begin{lemma}
\label{lem:cone-closed}
유한 개의 벡터가 생성하는 원뿔은 닫힌 볼록 집합이다.
\end{lemma}

\begin{proof}
볼록인 것은 정의에서 바로 나온다. 닫혀 있음을 보이려고 먼저 $K$의 모든 점이 일차 독립인 생성 벡터만으로 음이 아닌 결합으로 쓰임을 보인다(Carathéodory). $y=\sum_{j\in S}\mu_ja_j$, $\mu_j>0$에서 $\{a_j\}_{j\in S}$가 일차 종속이면 $\sum_{j\in S}\gamma_ja_j=0$인 $\gamma\ne0$이 있고, 부호를 바꾸어 어떤 $\gamma_j>0$이라 해도 된다. $t=\min_{\gamma_j>0}\mu_j/\gamma_j$로 두면 $\mu_j-t\gamma_j\ge0$이고 적어도 하나가 $0$이 되므로, $y$를 생성 벡터가 하나 적은 결합으로 쓸 수 있다. 이를 되풀이하면 된다. 그러므로 $K$는 일차 독립인 부분집합 $S$들이 생성하는 원뿔 $K_S$의 유한 합집합이다. 각 $K_S$는 일차 독립인 열을 가진 행렬 $A_S$에 대해 $\{A_S\mu:\mu\ge0\}$인데, $A_S$가 단사이고 그 상이 닫힌 부분공간이므로 $A_S$는 $\mathbb R^{|S|}$와 그 상 사이의 위상동형이고, 닫힌 집합 $\{\mu\ge0\}$의 상은 닫혀 있다. 닫힌 집합의 유한 합집합은 닫혀 있다.
\end{proof}

\begin{theorem}[Farkas]
\label{thm:farkas}
$A\in\mathbb R^{n\times k}$, $b\in\mathbb R^n$에 대해 다음 둘 가운데 정확히 하나가 성립한다.
(i) $A\mu=b$, $\mu\ge0$인 $\mu$가 있다.
(ii) $A^{\mathsf T}d\ge0$, $b^{\mathsf T}d<0$인 $d$가 있다.
\end{theorem}

\begin{proof}
둘이 함께 성립하면 $0>b^{\mathsf T}d=\mu^{\mathsf T}A^{\mathsf T}d\ge0$이어서 모순이다. (i)이 성립하지 않는다고 하자. $A$의 열이 생성하는 원뿔 $K$는 \cref{lem:cone-closed}로 닫힌 볼록 집합이고 $b\notin K$이다. 42장의 사영 정리로 $b$의 사영 $p=P_K(b)$가 있고 모든 $y\in K$에 대해 $(b-p)^{\mathsf T}(y-p)\le0$이다. $d=p-b\ne0$으로 두자. $y=2p$와 $y=0$을 넣으면 $d^{\mathsf T}p=0$이다. 그러면 모든 $y\in K$에서 $d^{\mathsf T}y\ge d^{\mathsf T}p=0$이고, 특히 $A$의 열에 대해 $A^{\mathsf T}d\ge0$이다. 또 $b^{\mathsf T}d=(p-d)^{\mathsf T}d=-\|d\|^2<0$이다.
\end{proof}

기하적으로 (ii)는 $b$와 원뿔 $K$를 가르는 초평면이 원점을 지난다는 뜻이다. Farkas가 1902년에 증명한 이 보조정리\cite{farkas1902}는 선형 부등식의 이론 전체의 뿌리이다.

\subsection{KKT 조건}

\begin{theorem}[Karush--Kuhn--Tucker, 선형 제약]
\label{thm:kkt-linear}
제약이 모두 일차식 $c_i(x)=a_i^{\mathsf T}x-b_i$이고 $f$가 연속 미분 가능하다고 하자. $x_\ast$가 국소 최소점이면 다음을 만족하는 $\lambda^\ast$가 있다.
\begin{align*}
&\nabla f(x_\ast)=\sum_{i\in\mathcal E\cup\mathcal I}\lambda_i^\ast a_i,&&c_i(x_\ast)=0\ (i\in\mathcal E),\quad c_i(x_\ast)\ge0\ (i\in\mathcal I),\\
&\lambda_i^\ast\ge0\ (i\in\mathcal I),&&\lambda_i^\ast c_i(x_\ast)=0\ (i\in\mathcal E\cup\mathcal I).
\end{align*}
\end{theorem}

\begin{proof}
$x_\ast$에서 움직일 수 있는 방향은 $a_i^{\mathsf T}d=0$($i\in\mathcal E$), $a_i^{\mathsf T}d\ge0$($i\in\mathcal A(x_\ast)\cap\mathcal I$)인 $d$이다. 그런 $d$와 작은 $t>0$에 대해 $x_\ast+td$는 활성 제약을 지키고, 활성이 아닌 부등식은 연속성으로 작은 $t$에서 여전히 만족되므로 가능해이다. 국소 최소이므로 $f(x_\ast+td)\ge f(x_\ast)$이고, $t\downarrow0$으로 보내면 $\nabla f(x_\ast)^{\mathsf T}d\ge0$이다. 등식 $a_i^{\mathsf T}d=0$을 두 부등식 $\pm a_i^{\mathsf T}d\ge0$으로 쓰면, 이것은 \cref{thm:farkas}의 (ii)가 행렬 $A=[a_i,-a_i\ (i\in\mathcal E);\ a_i\ (i\in\mathcal A\cap\mathcal I)]$와 $b=\nabla f(x_\ast)$에 대해 성립하지 않는다는 뜻이다. 따라서 (i)이 성립해 $\nabla f(x_\ast)$가 이 열들의 음이 아닌 결합이다. 등식 제약의 두 계수의 차를 $\lambda_i^\ast$로, 활성 부등식의 계수를 $\lambda_i^\ast\ge0$으로, 활성이 아닌 부등식의 승수를 $0$으로 두면 된다.
\end{proof}

마지막 조건 $\lambda_i^\ast c_i(x_\ast)=0$을 \term{상보 여유}{complementary slackness} 조건이라 한다. 활성이 아닌 제약의 승수는 $0$이다. 승수 $\lambda_i^\ast$는 제약 $c_i\ge0$을 조금 늦출 때 최솟값이 바뀌는 비율, 곧 그 제약의 그림자 가격이다. 비선형 제약에서는 가능한 방향의 집합을 제약의 선형화로 정확히 나타낼 수 있다는 가정이 더 필요하다. 활성 제약의 기울기가 일차 독립이라는 \term{LICQ}{linear independence constraint qualification}가 그 대표적인 조건이며, 그 아래서 같은 결론이 성립한다\cite{nocedal2006}. 이 조건들은 Karush가 1939년의 석사 논문에서\cite{karush1939}, Kuhn과 Tucker가 1951년에 독립적으로 얻었다\cite{kuhn1951}.

\section{라그랑주 쌍대성}
\label{sec:lagrange-duality}

부등식 제약만 있는 문제 $\min f(x)$, $c(x)\ge0$에 대해 \term{쌍대 함수}{dual function}
\[
q(\lambda)=\inf_x\mathcal L(x,\lambda)=\inf_x\bigl(f(x)-\lambda^{\mathsf T}c(x)\bigr),\qquad\lambda\ge0
\]
를 정의하고, \term{쌍대 문제}{dual problem} $\max_{\lambda\ge0}q(\lambda)$를 생각한다. $q$는 $\lambda$의 일차 함수들의 하한이므로 언제나 오목이다. 원래 문제가 볼록이 아니어도 쌍대 문제는 볼록 문제이다.

\begin{proposition}[약한 쌍대성]
가능해 $x$와 $\lambda\ge0$에 대해 $q(\lambda)\le f(x)$이다.
\end{proposition}

\begin{proof}
$q(\lambda)\le f(x)-\lambda^{\mathsf T}c(x)\le f(x)$이다. 앞의 부등식은 하한의 정의이고, 뒤의 것은 $\lambda\ge0$, $c(x)\ge0$에서 나온다.
\end{proof}

그래서 쌍대 문제의 최댓값은 원래 문제의 최솟값의 하한이다. 둘의 차를 \term{쌍대 간격}{duality gap}이라 한다. 쌍대 간격은 반복 알고리즘의 멈춤 판정으로 쓰인다. 지금의 가능해와 승수에서 계산한 $f(x)-q(\lambda)$가 작으면, 최솟값을 모르면서도 지금의 해가 그만큼 이내로 최적임을 안다. 64장의 lasso가 이것을 쓴다. $f$와 $-c_i$가 볼록이고 엄격한 가능해($c(x)>0$인 $x$)가 있으면 쌍대 간격이 $0$이라는 것이 Slater 조건 아래의 \term{강한 쌍대성}{strong duality}이다\cite{boyd2004convex}. 선형계획에서는 이 조건조차 필요 없다.

\section{선형계획}
\label{sec:linear-programming}

\subsection{쌍대 정리}

\term{선형계획}{linear programming}의 표준형은
\[
\text{(P)}\quad\min_xc^{\mathsf T}x\quad\text{조건: }Ax=b,\ x\ge0
\]
이고 $A\in\mathbb R^{m\times n}$이다. 라그랑주 함수 $c^{\mathsf T}x-y^{\mathsf T}(Ax-b)-s^{\mathsf T}x$의 하한은 $c-A^{\mathsf T}y-s=0$일 때만 유한하므로, 쌍대 문제는
\[
\text{(D)}\quad\max_{y,s}b^{\mathsf T}y\quad\text{조건: }A^{\mathsf T}y+s=c,\ s\ge0
\]
이다. 약한 쌍대성은 직접 보인다. $c^{\mathsf T}x-b^{\mathsf T}y=(A^{\mathsf T}y+s)^{\mathsf T}x-(Ax)^{\mathsf T}y=s^{\mathsf T}x\ge0$이다.

\begin{theorem}[선형계획의 쌍대 정리]
\label{thm:lp-duality}
(P)에 최적해가 있으면 (D)에도 최적해가 있고 두 최적값이 같다. 이때 $x$와 $(y,s)$가 각각 최적인 것은 $x_is_i=0$($i=1,\dots,n$)인 것과 동치이다.
\end{theorem}

\begin{proof}
최적값을 $p^\ast$라 하자. (P)의 최적해 $x_\ast$는 일차식 제약 아래의 국소 최소점이므로 \cref{thm:kkt-linear}에서 $c=A^{\mathsf T}y+s$, $s\ge0$, $s_ix_{\ast i}=0$인 승수 $y$, $s$가 있다. 그러면 $(y,s)$는 (D)의 가능해이고 $b^{\mathsf T}y=c^{\mathsf T}x_\ast-s^{\mathsf T}x_\ast=p^\ast$이다. 약한 쌍대성으로 (D)의 어느 가능해도 $p^\ast$를 넘지 못하므로 $(y,s)$가 (D)의 최적해이다. 마지막 주장은 $c^{\mathsf T}x-b^{\mathsf T}y=s^{\mathsf T}x=\sum x_is_i$이고 각 항이 음이 아니라는 데서 나온다.
\end{proof}

\subsection{심플렉스 방법}

(P)의 가능 영역은 다면체이고, 일차 함수의 최솟값은 (있다면) 꼭짓점에서 얻어진다. $A$의 열 $m$개를 골라 정칙인 $B$를 만들고 나머지 변수를 $0$으로 둔 해 $x_B=B^{-1}b$가 $x_B\ge0$이면 \term{기저 가능해}{basic feasible solution}라 하는데, 이것이 정확히 다면체의 꼭짓점이다. Dantzig의 \term{심플렉스 방법}{simplex method}은 꼭짓점에서 이웃한 꼭짓점으로 목적 함수가 줄어드는 쪽으로 옮겨 간다\cite{dantzig1963}. 비기저 변수 $j$의 축소 비용 $\bar c_j=c_j-c_B^{\mathsf T}B^{-1}a_j$가 음수이면 그 변수를 $0$에서 늘려 목적 함수를 줄일 수 있다. 그 변수를 기저에 들이고, 늘리다가 가장 먼저 $0$이 되는 기저 변수를 내보낸다(비율 검사). 모든 $\bar c_j\ge0$이면 $y=B^{-\mathsf T}c_B$가 쌍대 가능해이고 상보 여유가 성립하므로 \cref{thm:lp-duality}로 최적이다.

들일 변수를 고르는 규칙은 여럿이다. Dantzig의 원래 규칙은 $\bar c_j$가 가장 음수인 것을 고른다. 기저 변수가 $0$인 퇴화한 꼭짓점에서는 피벗해도 목적 함수가 줄지 않아 같은 기저들을 끝없이 돌 수 있는데, Bland는 들일 변수와 내보낼 변수 모두 번호가 가장 작은 것을 고르면 순환하지 않음을 보였다\cite{bland1977}. 기저가 유한하므로 그러면 방법은 유한 번에 끝난다.

심플렉스 방법은 실제로는 대개 $m$의 몇 배 정도의 피벗으로 끝나지만, 최악의 경우는 지수적이다. Klee와 Minty는 1972년에 $n$차원 입방체를 찌그러뜨린 문제
\[
\max\sum_{j=1}^n2^{n-j}x_j\quad\text{조건: }2\sum_{j<i}2^{i-j}x_j+x_i\le5^i\ (i=1,\dots,n),\quad x\ge0
\]
에서 Dantzig의 규칙이 입방체의 꼭짓점 $2^n$개를 모두 지나는 것을 보였다\cite{klee1972}. 유리수로 정확히 계산해 보면 $n=2,\dots,10$에서 피벗 수가 정확히 $3,7,15,\dots,1023=2^n-1$이다. Bland의 규칙도 이 문제에서는 $3,5,9,15,25,41,67,109,177$로, 대략 $\varphi\approx1.62$배씩 늘어 역시 지수적이다.

\subsection{내부점법}

Karmarkar는 1984년에 다항 시간이면서 실제로도 빠른 \term{내부점법}{interior-point method}을 내놓았다\cite{karmarkar1984}. 오늘날 표준은 원-쌍대 내부점법이다\cite{wright1997pdipm}. \cref{thm:lp-duality}의 최적성 조건
\[
Ax=b,\qquad A^{\mathsf T}y+s=c,\qquad x_is_i=0,\qquad x,s\ge0
\]
에서 상보 조건을 $x_is_i=\tau$로 늦춘 계는 $\tau>0$마다 유일한 해 $(x(\tau),y(\tau),s(\tau))$를 가지며(원 문제와 쌍대 문제에 모두 $x>0$, $s>0$인 가능해가 있고 $A$의 행이 일차 독립일 때), $\tau\downarrow0$에 따라 이 해들이 그리는 곡선을 \term{중심 경로}{central path}라 한다. 내부점법은 $x,s>0$을 지키며 이 비선형 방정식계에 34장의 Newton 방법을 한 걸음씩 쓰고, 걸음마다 $\tau$를 줄인다. 이론적으로는 $O(\sqrt n\log(1/\epsilon))$번의 반복으로 정확도 $\epsilon$에 이르고, Mehrotra의 예측--교정 방법\cite{mehrotra1992}을 쓰면 실제로는 문제의 크기와 거의 상관없이 수십 번이면 된다. 반복마다 $m\times m$ 행렬 $AD A^{\mathsf T}$($D$는 대각)의 Cholesky 분해가 하나 드는 것이 비용이다.

무작위로 만든 $m\times n$ 문제를 HiGHS(scipy 1.18.1의 \code{linprog})로 풀면 $(m,n)=(50,100)$, $(200,400)$, $(1000,2000)$에서 내부점법의 반복이 $12$, $16$, $21$번이고 쌍대 심플렉스의 피벗이 $68$, $348$, $2{,}627$번이다. 내부점법의 반복은 비싸지만 수가 거의 늘지 않는다. 66장의 분위 회귀는 선형계획이며, stardust는 그것을 Portnoy와 Koenker의 Frisch--Newton 내부점법으로 풀고\cite{portnoy1997} 마지막에 정확한 꼭짓점 해를 찾는다.

\section{상자 제약 최소제곱}
\label{sec:bounded-lsq}

\subsection{음이 아닌 최소제곱}

통계에서 자주 만나는 제약 문제는 선형 최소제곱에 부호나 범위의 제약을 단 것이다. 농도, 강도, 혼합 비율이 그렇다.
\[
\min_x\frac12\|Ax-b\|^2\quad\text{조건: }x\ge0
\]
을 \term{음이 아닌 최소제곱}{nonnegative least squares, NNLS}이라 한다. 목적 함수가 볼록이고 제약이 일차식이므로 \cref{thm:kkt-linear}의 조건이 최적의 필요충분조건이다(볼록 문제에서는 KKT 점이 전역 최소이다). $w=A^{\mathsf T}(b-Ax)$, 곧 음의 기울기로 쓰면 조건은
\[
x\ge0,\qquad w\le0,\qquad x_iw_i=0
\]
이다. 양수인 성분에서는 기울기가 $0$이고, $0$인 성분에서는 기울기가 그 성분을 늘리는 쪽으로 내리막이 아니다.

Lawson과 Hanson의 알고리즘\cite{lawson1995}은 양수일 수 있는 변수의 집합 $P$(수동 집합)를 키워 가는 활성 집합 방법이다(\cref{alg:nnls}). 처음에는 $P=\varnothing$, $x=0$이다. $P$ 밖에서 $w_j$가 가장 큰 변수가 양수이면 그것을 $P$에 들이고, $P$의 열만으로 제약 없는 최소제곱 $z$를 푼다. $z$가 $P$에서 모두 양수이면 $x=z$로 받아들인다. 아니면 $x$에서 $z$ 쪽으로, 처음 $0$에 닿는 변수가 생길 때까지만 가고 그 변수를 $P$에서 뺀 뒤 다시 푼다.

\begin{algorithm}[t]
\begin{algorithmic}[1]
\Require 행렬 $A$, 벡터 $b$
\Ensure $\min\|Ax-b\|$, $x\ge0$의 해
\State $P\gets\varnothing$, $x\gets0$, $w\gets A^{\mathsf T}b$
\While{$P$ 밖에 $w_j>0$인 $j$가 있다}
  \State $j\gets\arg\max_{j\notin P}w_j$, $P\gets P\cup\{j\}$
  \Loop
    \State $z_P\gets\arg\min\|A_Pz_P-b\|$, $z_j\gets0$ ($j\notin P$) \Comment{제약 없는 최소제곱}
    \If{모든 $i\in P$에서 $z_i>0$} \State \textbf{break} \EndIf
    \State $\alpha\gets\min_{i\in P,\,z_i\le0}x_i/(x_i-z_i)$, $x\gets x+\alpha(z-x)$
    \State $x_i=0$인 $i$를 $P$에서 뺀다
  \EndLoop
  \State $x\gets z$, $w\gets A^{\mathsf T}(b-Ax)$
\EndWhile
\end{algorithmic}
\caption{Lawson--Hanson의 NNLS}
\label{alg:nnls}
\end{algorithm}

바깥 반복마다 목적 함수가 엄격히 줄고 그때의 $x$는 수동 집합 $P$에 대한 제약 없는 최소제곱의 해이므로, 같은 $P$가 두 번 나올 수 없다. $P$의 경우의 수가 유한하므로 알고리즘은 유한 번에 끝나고, 끝날 때는 KKT 조건이 성립한다. 실제 반복 수는 대개 $n$ 안팎이다. 무작위 $20\times10$ 문제에서 제약 없는 최소제곱 해의 음수 성분은 $6$개이고, NNLS의 해는 양수 성분이 $5$개이다. 그 해에서 $w$는 양수 성분에서 크기가 $2.9\times10^{-15}$ 이하이고 $0$인 성분에서 가장 큰 값이 $-1.69$이므로 KKT 조건이 반올림 수준으로 성립한다. 제약 없는 해의 음수를 $0$으로 자른 것과 NNLS의 해는 일반적으로 다르다는 데 주의하자. 성분 하나를 $0$으로 묶으면 나머지 성분의 최적값이 바뀐다.

\subsection{상하한과 BVLS}

제약을 $l\le x\le u$로 넓힌 문제를 Stark와 Parker의 \term{BVLS}{bounded-variable least squares}가 같은 생각으로 푼다\cite{stark1995}. 변수는 자유, 아래 경계에 묶임, 위 경계에 묶임의 세 상태 가운데 하나이고, 묶인 변수의 기울기가 경계 안쪽을 가리키면 풀어 주고, 자유 변수의 최소제곱 해가 경계를 넘으면 그 경계로 묶는다. 47장의 반사 신뢰 영역 방법도 이 문제에 쓸 수 있다. 위의 $20\times10$ 문제에 $-0.2\le x\le0.2$를 주면 BVLS는 $5$번, 반사 신뢰 영역은 $11$번 반복해 같은 최솟값 $4.3763$에 이르고, 네 변수가 경계에 묶인다. 활성 집합 방법은 변수가 수백 개 이하일 때 빠르고 정확한 해를 주며, 반사 신뢰 영역은 큰 희소 문제로 넓힐 수 있다.

\begin{instardust}[\code{docs/design/optimize.md} §1.1, §3.4]
stardust의 \code{sd.optimize}에는 선형계획(\code{linprog})과 일반 제약 최적화(SLSQP, \code{trust-constr})가 아직 없고(설계 문서는 이것들을 나중에 따로 정하기로 미루었다), 제약은 상자 제약만 다룬다. 상자 제약 최소제곱은 \code{lsq\_linear}(\code{'trf'}와 \code{'bvls'})과 \code{nnls}이다(\code{csrc/optimize/linear.cpp}). \code{lsq\_linear}은 scipy처럼 먼저 LAPACK \code{gelsd}로 제약 없는 해를 구해, 그것이 경계 안이면 반복 없이 끝낸다. \code{'bvls'}가 돌려주는 $x$는 언제나 경계 안이도록, 경계에 묶인 변수는 정확히 경계 값으로 두고 반올림만큼 넘은 값은 자른다. \code{nnls}는 scipy 1.18.1의 C 판 Lawson--Hanson을 옮겼는데, 짧은 내적에서 Accelerate를 부르는 비용이 시간을 지배해 쌍대 값의 내적과 열 하나의 Householder 반사를 누산기 네 개의 고리로 바꾸었다. 그래서 $x$는 $0$인 자리가 scipy와 같고 나머지는 끝자리가 다르다. 설계 문서는 그 차이가 남은 열의 조건수에 비례해 조건이 좋은 문제에서 $9$ ULP, 열이 비슷한 지수 함수 적합에서 $652$ ULP까지라고 적는다. 합의 순서는 길이로만 정해지므로 실행마다, 스레드 수와 무관하게 같다.
\end{instardust}

\section{맺음}

제약 최적화의 이론은 Farkas 보조정리 하나에서 자란다. 그것이 KKT 조건을 주고, 승수는 쌍대 문제의 변수가 되며, 쌍대 문제는 원래 문제의 하한과 최적성의 증명서를 준다. 선형계획에서는 이 그림이 완전해서 원래 문제와 쌍대 문제의 최적값이 같고, 상보 여유 조건이 두 문제의 최적해를 잇는다. 심플렉스 방법은 꼭짓점을 따라 이 조건을 찾고, 내부점법은 중심 경로를 따라 다가간다. 최악의 경우 하나는 지수적이고 다른 하나는 다항적이지만, 실제로는 둘 다 쓰인다. 통계에서 가장 자주 만나는 제약은 상자 제약이고, NNLS와 BVLS의 활성 집합 방법은 KKT 조건을 하나씩 맞추어 가는 가장 직접적인 알고리즘이다. 이것으로 최적화의 부를 마치고, 다음 부에서는 확률과 통계적 추론으로 넘어간다.

볼록 최적화와 쌍대성은 Boyd와 Vandenberghe의 책\cite{boyd2004convex}, 제약 최적화의 알고리즘은 Nocedal과 Wright의 책\cite{nocedal2006} 12--19장, 내부점법은 Wright의 책\cite{wright1997pdipm}이 자세하다.

\chapterreferences
